% SGA 3, Expose IV, English reconstruction.
% Scope: all of Section 3, stopping immediately before the Section 4 heading.
% Authority: Polo--Gille Exp4-13oct24.pdf, component-local pages 10--20
% (combined-reader pages 198--208; printed pages 186--196).
% Jacob Reinhold's CC BY 4.0 Markdown lineage was used only as a phrasing
% comparison. The Polo--Gille PDF controls prose, formulas, numbering, and notes.

\setcounter{footnote}{9}

% Source page: Exp4 component-local p. 10; combined-reader p. 198; printed p. 186.
\SGARefTarget{sga3:iv:section:3}\section*{3. 普遍有效等价关系}
\addcontentsline{toc}{section}{3. 普遍有效等价关系}

\SGARefTarget{sga3:iv:subsection:3:1}\subsection*{3.1. 等价关系: 定义}
\addcontentsline{toc}{subsection}{\SGARefLink{sga3:iv:subsection:3:1}{3.1}. 等价关系: 定义}

\noindent\SGARefTarget{sga3:iv:definition:3:1:1}\textbf{定义 3.1.1.}
\textit{对象 \(X\in\operatorname{Ob}\mathscr C\) 中的一个
\(\mathscr C\)-等价关系, 是指函子 \(X\times X\) 的一个可表子函子
\(R\), 使得对每个 \(S\in\operatorname{Ob}\mathscr C\),
\(R(S)\) 是 \(X(S)\) 上某个等价关系的图像.}

\SGARefLink{sga3:iv:definition:3:1:1}{这一定义}特别适用于范畴
\(\widehat{\mathscr C}\). 若将 \(X\) 视为 \(\widehat{\mathscr C}\)
的对象, 则 \(X\) 中的一个 \(\widehat{\mathscr C}\)-等价关系,
就是 \(X\times X\) 的一个子函子 \(R\), 满足对
\(\mathscr C\) 的每个对象 \(S\), \(R(S)\) 都是 \(X(S)\)
上某个等价关系的图像.
\footnote{编者注: 这一条件显然是必要的. 反之, 若对每个
\(S\in\operatorname{Ob}\mathscr C\), \(R(S)\) 都是某个等价关系的图像,
则这一等价关系扩张为对每个 \(F\in\operatorname{Ob}\widehat{\mathscr C}\)
定义的 \(R(F)\): 规定两个态射 \(\varphi,\psi:F\to R\) 等价,
是指对每个 \(S\in\operatorname{Ob}\mathscr C\) 和每个 \(x\in F(S)\),
元素 \(\varphi(x)\) 与 \(\psi(x)\) 在 \(X(S)\) 中等价.}

若 \(R\) 是 \(X\) 中的一个 \(\mathscr C\)-等价关系,
用 \(p_i\) 表示投影 \(\operatorname{pr}_i:X\times X\to X\)
诱导的从 \(R\) 到 \(X\) 的态射. 于是有图
\[
  p_1,p_2:R\rightrightarrows X.
\]

\medskip
\noindent\SGARefTarget{sga3:iv:definition:3:1:2}\textbf{定义 3.1.2.}
\textit{若 \(up_1=up_2\), 则称态射 \(u:X\to Z\)
\emph{与 \(R\) 相容}. \(\mathscr C\) 中作为箭头对 \((p_1,p_2)\)
的余核的对象, 也称为 \(X\) 对 \(R\) 的\emph{商对象},
记为 \(X/R\).}

因而有正合图
\[
  R\overset{p_1}{\underset{p_2}{\rightrightarrows}}
  X\xrightarrow{p}X/R,
\]
且 \(X/R\) 表示以下协变函子
\[
  \operatorname{Hom}_{\mathscr C}(X/R,Z)
  =\{\text{从 \(X\) 到 \(Z\) 且与 \(R\) 相容的态射}\}.
\]
由于商对象在 \(\mathscr C\) 中取, 商 \(X/R\) 只要存在就是唯一的.

这些定义立即推广到 \(X\) 中的 \(\widehat{\mathscr C}\)-等价关系.
但应当注意, 将 \(\mathscr C\) 的对象等同于其在
\(\widehat{\mathscr C}\) 中的像, 与商的构造并不交换:
在 \(\mathscr C\) 中 \(X\) 对 \(R\) 的商 \(X/R\),
并\emph{不能预先认为}是在 \(\widehat{\mathscr C}\) 中
\(X\) 对 \(R\) 的商. 因此, 在涉及取商的问题中,
必须避免不加分辨地将 \(\mathscr C\) 等同于其在
\(\widehat{\mathscr C}\) 中的像.
\footnote{编者注: 以下例子也预示了本篇的后续内容.
设 \(G\) 为 \(\mathscr C\)-群, \(H\) 为一个
\(\mathscr C\)-子群, \(R\) 为 \(G\) 中由
\(G\times H\to G\times G\), \((g,h)\mapsto(g,gh)\)
定义的等价关系 (参见 \SGARefLink{sga3:iv:subsection:3:2}{3.2}).
由 \(Q(S)=G(S)/H(S)\) 定义的函子 \(Q\) 是
\(\widehat{\mathscr C}\) 中的商
(将 \SGARefLink{sga3:iv:proposition:4:4:9}{4.4.9} 应用于最粗拓扑;
参见 \SGARefLink{sga3:iv:scholium:4:4:2}{4.4.2}),
但一般不是所希望得到的商. 例如, 当 \(\mathscr C=\mathbf{Sch}\) 时,
有仿射群概形的正合列
\[
  1\longrightarrow\mu_2\longrightarrow\mathbf G_m
  \xrightarrow{p}\mathbf G_m\longrightarrow1
\]
将 \(\mathbf G_m\) 等同于商 \(\mathbf G_m/\mu_2\).
此外, 由于 \(p\) 是有限且局部自由的, 在 fppf 拓扑的层构成的
更大范畴中, \(\mathbf G_m\) 是 \(\mathbf G_m\) 对 \(\mu_2\)
的层商; 参见 \SGARefLink{sga3:iv:theorem:4:6:6:item:ii:01}{4.6.6(ii)}
和 \SGARefLink{sga3:iv:corollary:6:3:2}{6.3.2}.
而 \(\widehat{\mathscr C}\) 中的商 \(Q\) 不同构于 \(\mathbf G_m\),
因为例如 \(Q(\mathbf Z)=\{1\}\), 而
\(\mathbf G_m(\mathbf Z)=\{\pm1\}\). 因此 \(Q\) 不是 fppf 层,
\emph{更不可能}可表.}

% Source page: Exp4 component-local p. 11; combined-reader p. 199; printed p. 187.
在下文中, 简称\emph{等价关系}时, 指的是
\(\widehat{\mathscr C}\)-等价关系; 必要时, 将明确说明所指的是
\(\mathscr C\)-等价关系.
\footnote{编者注: 即使是 \(X\) 中的 \(\widehat{\mathscr C}\)-等价关系,
我们关心的问题也是 \(\mathscr C\) 中的商是否存在.}

\medskip
\noindent\SGARefTarget{sga3:iv:definition:3:1:3}\textbf{定义 3.1.3.}
\textit{若 \(X\) 是 \(\mathscr C\) 中 \(S\) 上的对象,
则 \emph{\(S\) 上 \(X\) 中的等价关系} 是指 \(X\) 中的一个等价关系
\(R\), 使结构态射 \(X\to S\) 与 \(R\) 相容.}

此时, 典范单态射 \(R\to X\times X\) 经过单态射
\[
  X\times_S X\longrightarrow X\times X
\]
分解, 并定义 \(\mathscr C_{/S}\) 的对象 \(X\to S\) 中的一个等价关系.
商 \(X/R\) 只要存在, 就带有到 \(S\) 的典范态射,
而 \(\mathscr C_{/S}\) 中的相应对象, 就是
\(X\in\operatorname{Ob}\mathscr C_{/S}\) 对上述等价关系的商.
反之, 若 \(S\) 是 \(\mathscr C\) 的可平方对象,
且 \(Y\to S\) 是 \(\mathscr C_{/S}\) 中 \(X\) 对此等价关系的商,
则 \(Y\) 是 \(\mathscr C\) 中 \(X\) 对 \(R\) 的商.
无论如何, 我们不必考虑那些在 \(\mathscr C_{/S}\) 中是商,
却尚不是 \(\mathscr C\) 中的商的情形.

\medskip
\noindent\SGARefTarget{sga3:iv:definition:3:1:4}\textbf{定义 3.1.4.}
\textit{设 \(X\) (分别地, \(X'\)) 是 \(\mathscr C\) 中带有等价关系
\(R\) (分别地, \(R'\)) 的对象. 若以下等价条件成立,
则称态射
\[
  u:X\longrightarrow X'
\]
\emph{与 \(R\) 和 \(R'\) 相容}:}
\begin{enumerate}
  \item[(i)]\SGARefTarget{sga3:iv:definition:3:1:4:item:i:01}
  \textit{对每个 \(S\in\operatorname{Ob}\mathscr C\),
  \(X(S)\) 中模 \(R(S)\) 同余的两点, 被 \(u\) 映到
  \(X'(S)\) 中模 \(R'(S)\) 同余的两点;}
  \item[(ii)]\SGARefTarget{sga3:iv:definition:3:1:4:item:ii:01}
  \textit{存在一个必然唯一的态射 \(R\to R'\), 使图}
  \[
    \begin{array}{ccc}
      R & \longrightarrow & R'\\
      \big\downarrow && \big\downarrow\\
      X\times X & \xrightarrow{u\times u} & X'\times X'
    \end{array}
  \]
  \textit{交换.}
\end{enumerate}

% Source page: Exp4 component-local p. 12; combined-reader p. 200; printed p. 188.
由 \(X/R\) 的泛性质, 只要 \(X/R\) 和 \(X'/R'\) 存在,
就有唯一的态射 \(v\) 使图
\[
  \begin{array}{ccc}
    X & \xrightarrow{p} & X/R\\
    {\scriptstyle u}\big\downarrow && \big\downarrow{\scriptstyle v}\\
    X' & \xrightarrow{p'} & X'/R'
  \end{array}
\]
交换.

\medskip
\noindent\SGARefTarget{sga3:iv:definition:3:1:5}\textbf{定义 3.1.5.}
\textit{若以下等价条件成立, 则称 \(X\) 的子对象
(即一个可表子函子) \(Y\) \emph{在等价关系 \(R\) 下稳定}:}
\begin{enumerate}
  \item[(i)]\SGARefTarget{sga3:iv:definition:3:1:5:item:i:01}
  \textit{对每个 \(S\in\operatorname{Ob}\mathscr C\),
  \(X(S)\) 的子集 \(Y(S)\) 在 \(R(S)\) 下稳定;}
  \item[(ii)]\SGARefTarget{sga3:iv:definition:3:1:5:item:ii:01}
  \textit{\(Y\) 在 \(p_1\) 和 \(p_2\) 下的逆像这两个
  \(R\) 的子对象相同.}
\end{enumerate}

\(X\) 的稳定子对象的一个重要特例如下: 商 \(X/R\) 存在,
且 \(Y\) 是 \(X/R\) 的某个子对象在 \(X\) 中的逆像.

\medskip
\noindent\SGARefTarget{sga3:iv:definition:3:1:6}\textbf{定义 3.1.6.}
\textit{设 \(R\) 是 \(X\) 中的等价关系, \(X'\to X\) 是态射.
从 Cartesian 图
\[
  \begin{array}{ccc}
    R' & \longrightarrow & R\\
    \big\downarrow && \big\downarrow\\
    X'\times X' & \longrightarrow & X\times X
  \end{array}
\]
得到的 \(X'\) 中的等价关系 \(R'\), 称为 \(X\) 中的等价关系
\(R\) 在 \(X'\) 中的\emph{逆像}等价关系.
特别地, 若 \(X'\) 是 \(X\) 的子对象, 则称之为 \(R\) 在
\(X'\) 上\emph{诱导}的等价关系, 并记为 \(R_{X'}\).}

态射 \(X'\to X\) 与 \(R'\) 和 \(R\) 相容; 因而, 当商存在时,
它诱导态射 \(X'/R'\to X/R\)
(\SGARefLink{sga3:iv:definition:3:1:4}{3.1.4}).
若 \(X'\) 是 \(X\) 的子对象, 后文将会看到,
在某些情形下可以证明 \(X'/R'\to X/R\) 为单态射,
因而将 \(X'/R'\) 等同于 \(X/R\) 的一个子对象.
在这种情形下, 此子对象在 \(X\) 中的逆像是一个包含 \(X'\)
且在 \(R\) 下稳定的子对象: 这就是 \(X'\) 关于等价关系
\(R\) 的\emph{饱和化 (saturation)}.

\medskip
\noindent\SGARefTarget{sga3:iv:proposition:3:1:7}\textbf{命题 3.1.7.}
\textit{若 \(X\) 的子对象 \(Y\) 在 \(R\) 下稳定,
则对 \(i=1,2\), 以下两个方形都是 Cartesian 方形:}
\[
  \begin{array}{ccc}
    R_Y & \longrightarrow & R\\
    {\scriptstyle p_i}\big\downarrow && \big\downarrow{\scriptstyle p_i}\\
    Y & \longrightarrow & X.
  \end{array}
\]

证明是直接的.

% Source page: Exp4 component-local p. 13; combined-reader p. 201; printed p. 189.
\SGARefTarget{sga3:iv:subsection:3:2}\subsection*{3.2. 群的自由作用定义的等价关系}
\addcontentsline{toc}{subsection}{\SGARefLink{sga3:iv:subsection:3:2}{3.2}. 群的自由作用定义的等价关系}

\noindent\SGARefTarget{sga3:iv:definition:3:2:1}\textbf{定义 3.2.1.}
\textit{设 \(X\) 是 \(\mathscr C\) 的对象, \(H\) 是作用在 \(X\) 上的
\(\mathscr C\)-群, 即带有一个 \(\widehat{\mathscr C}\)-群态射}
\[
  H\longrightarrow\operatorname{Aut}(X).
\]
\textit{若以下等价条件成立, 则称群 \(H\) 在 \(X\) 上\emph{自由}作用:}
\begin{enumerate}
  \item[(i)]\SGARefTarget{sga3:iv:definition:3:2:1:item:i:01}
  \textit{对每个 \(S\in\operatorname{Ob}\mathscr C\),
  群 \(H(S)\) 在 \(X(S)\) 上自由作用;}
  \item[(ii)]\SGARefTarget{sga3:iv:definition:3:2:1:item:ii:01}
  \textit{在集合上由 \((h,x)\mapsto(hx,x)\) 定义的函子态射}
  \[
    H\times X\longrightarrow X\times X
  \]
  \textit{是单态射.}
\end{enumerate}

在 \SGARefLink{sga3:iv:definition:3:2:1}{前一定义} 中,
假设群 \(H\) 在 \(X\) 上左作用. 当然, 对群右作用的情形也有类似概念,
即给定从反群 \(H^{\circ}\) 到 \(\operatorname{Aut}(X)\) 的群态射.

若 \(H\) 在 \(X\) 上自由作用, 则 \(H\times X\) 在
\SGARefLink{sga3:iv:definition:3:2:1:item:ii:01}{(ii)} 中的态射下的像,
是 \(X\) 中的一个等价关系, 称为\emph{由 \(H\) 在 \(X\) 上的作用
定义的等价关系}. 当 \(X\) 对此等价关系的商存在时,
记为 \(H\backslash X\) (若 \(H\) 右作用, 则记为 \(X/H\)).
它表示以下协变函子: 对 \(\mathscr C\) 的每个对象 \(Z\),
\[
  \operatorname{Hom}(H\backslash X,Z)
  =\{\text{从 \(X\) 到 \(Z\) 且在 \(H\) 下不变的态射}\},
\]
其中, 称态射 \(f:X\to Z\) 在 \(H\) 下不变, 是指对每个
\(S\in\operatorname{Ob}\mathscr C\), 相应的态射
\(X(S)\to Z(S)\) 在群 \(H(S)\) 下不变.

\medskip
\noindent\SGARefTarget{sga3:iv:lemma:3:2:2}\textbf{引理 3.2.2.}
\textit{在 \SGARefLink{sga3:iv:definition:3:2:1}{3.2.1} 的假设下,
设 \(Y\) 是 \(X\) 的子对象. 以下条件等价:}
\begin{enumerate}
  \item[(i)]\SGARefTarget{sga3:iv:lemma:3:2:2:item:i:01}
  \textit{\(Y\) 在 \(H\) 定义的等价关系下稳定
  (\SGARefLink{sga3:iv:definition:3:1:5}{3.1.5});}
  \item[(ii)]\SGARefTarget{sga3:iv:lemma:3:2:2:item:ii:01}
  \textit{对每个 \(S\in\operatorname{Ob}\mathscr C\),
  \(X(S)\) 的子集 \(Y(S)\) 在 \(H(S)\) 下稳定;}
  \item[(iii)]\SGARefTarget{sga3:iv:lemma:3:2:2:item:iii:01}
  \textit{存在一个必然唯一的态射 \(f\), 使图}
  \[
    \begin{array}{ccc}
      H\times Y & \xrightarrow{f} & Y\\
      \big\downarrow && \big\downarrow\\
      H\times X & \longrightarrow & X
    \end{array}
  \]
  \textit{交换.}
\end{enumerate}
\textit{在这些条件下, \(f\) 定义一个 \(\widehat{\mathscr C}\)-群态射}
\[
  H\longrightarrow\operatorname{Aut}(Y),
\]
\textit{并且, \(H\) 的这一作用在 \(Y\) 中定义的等价关系,
恰好是 \(H\) 的作用在 \(X\) 中定义的等价关系在 \(Y\) 上诱导的等价关系.}

% Source page: Exp4 component-local p. 14; combined-reader p. 202; printed p. 190.
证明是直接的. 对右作用当然也有类似陈述.
上面定义的 \(H\) 在 \(Y\) 上的作用, 称为给定的 \(H\) 在 \(X\)
上的作用在 \(Y\) 上诱导的作用.

现考虑两个 \(\mathscr C\)-群 \(H\) 和 \(G\), 以及群态射
\[
  u:H\longrightarrow G.
\]
则 \(H\) 以平移的方式作用在 \(G\) 上 (在集合上,
\(hg=u(h)g\)), 且这一作用自由当且仅当 \(u\) 为单态射.
当 \(G\) 对 \(H\) 的这一作用的商存在时, 记为 \(H\backslash G\).
类似地可定义 \(H\) 在 \(G\) 上的右作用和商 \(G/H\).
这些商关于所涉及的群具有函子性. 更确切地说,
对右商有以下 \SGARefLink{sga3:iv:lemma:3:2:3}{引理}.

\medskip
\noindent\SGARefTarget{sga3:iv:lemma:3:2:3}\textbf{引理 3.2.3.}
\textit{设 \(u:H\to G\) 和 \(u':H'\to G'\) 是
\(\mathscr C\)-群的单态射, 并假设给定 \(\mathscr C\)-群态射
\[
  f:G\longrightarrow G'
\]
则以下条件等价:}
\begin{enumerate}
  \item[(i)]\SGARefTarget{sga3:iv:lemma:3:2:3:item:i:01}
  \textit{\(f\) 与 \(H\) 和 \(H'\) 在 \(G\) 和 \(G'\)
  中定义的等价关系相容;}
  \item[(ii)]\SGARefTarget{sga3:iv:lemma:3:2:3:item:ii:01}
  \textit{对每个 \(S\in\operatorname{Ob}\mathscr C\), 有
  \(fu(H(S))\subseteq u'(H'(S))\);}
  \item[(iii)]\SGARefTarget{sga3:iv:lemma:3:2:3:item:iii:01}
  \textit{存在一个必然唯一的保乘法态射 \(g:H\to H'\), 使图}
  \[
    \begin{array}{ccc}
      H & \xrightarrow{g} & H'\\
      {\scriptstyle u}\big\downarrow && \big\downarrow{\scriptstyle u'}\\
      G & \xrightarrow{f} & G'
    \end{array}
  \]
  \textit{交换.}
\end{enumerate}
\textit{在这些条件下, 若商 \(G/H\) 和 \(G'/H'\) 存在,
则有唯一态射 \(\overline f\) 使图}
\[
  \begin{array}{ccc}
    G & \xrightarrow{f} & G'\\
    {\scriptstyle p}\big\downarrow && \big\downarrow{\scriptstyle p'}\\
    G/H & \xrightarrow{\overline f} & G'/H'
  \end{array}
\]
\textit{交换.}

第一个断言可化归到集合的情形来证明. 第二个立即由
\SGARefLink{sga3:iv:lemma:3:2:3:item:i:01}{(i)} 得到.

前面对于一般等价关系引入的概念, 都可以改述为当前情形中的概念.
这里只记下 \SGARefLink{sga3:iv:lemma:3:2:4}{如下引理},
其证明化归到集合的情形后是直接的.

% Source page: Exp4 component-local p. 15; combined-reader p. 203; printed p. 191.
\medskip
\noindent\SGARefTarget{sga3:iv:lemma:3:2:4}\textbf{引理 3.2.4.}
\textit{设 \(u:H\to G\) 是 \(\mathscr C\)-群的单态射,
\(G'\) 是 \(G\) 的 \(\mathscr C\)-子群.
\(G\) 的子对象 \(G'\) 在 \(H\) 定义的等价关系下稳定,
当且仅当 \(u\) 经过典范单态射 \(G'\to G\) 分解.
在 \SGARefLink{sga3:iv:lemma:3:2:4}{此条件} 下,
由 \(u\) 定义的 \(G\) 上的作用在 \(G'\) 上诱导的 \(H\) 的作用,
恰好是 \(u\) 所经过的单态射 \(H\to G'\) 产生的作用.}

\SGARefTarget{sga3:iv:subsection:3:3}\subsection*{3.3. 普遍有效等价关系}
\addcontentsline{toc}{subsection}{\SGARefLink{sga3:iv:subsection:3:3}{3.3}. 普遍有效等价关系}

\noindent\SGARefTarget{sga3:iv:definition:3:3:1}\textbf{定义 3.3.1.}
\textit{设 \(f:X\to Y\) 是态射. \emph{\(f\) 在 \(X\) 中定义的等价关系},
记为 \(\mathscr R(f)\), 是 \(X\) 中作为典范单态射}
\[
  X\times_Y X\longrightarrow X\times X.
\]
\textit{的像的 \(\widehat{\mathscr C}\)-等价关系.}

\medskip
\noindent\SGARefTarget{sga3:iv:definition:3:3:2}\textbf{定义 3.3.2.}
\textit{若下列条件成立, 则称 \(X\) 中的等价关系 \(R\) 为\emph{有效的}:}
\begin{enumerate}
  \item[(i)]\SGARefTarget{sga3:iv:definition:3:3:2:item:i:01}
  \textit{\(R\) 可表, 即它是一个 \(\mathscr C\)-等价关系;}
  \item[(ii)]\SGARefTarget{sga3:iv:definition:3:3:2:item:ii:01}
  \textit{商 \(Y=X/R\) 在 \(\mathscr C\) 中存在;}
  \footnote{编者注: ``在 \(\mathscr C\) 中''这几个字为新增内容.}
  \item[(iii)]\SGARefTarget{sga3:iv:definition:3:3:2:item:iii:01}
  \textit{图}
  \[
    R\overset{p_1}{\underset{p_2}{\rightrightarrows}}
    X\xrightarrow{p}Y
  \]
  \textit{使 \(R\) 成为 \(X\) 在 \(Y\) 上的纤维平方,
  即 \(R\) 是 \(p\) 定义的等价关系.}
\end{enumerate}

\medskip
\noindent\SGARefTarget{sga3:iv:scholium:3:3:2:1}\textbf{附论 3.3.2.1.}
\footnote{编者注: 为便于后文引用, 添加了编号
\SGARefLink{sga3:iv:scholium:3:3:2:1}{3.3.2.1}
(分别地, \SGARefLink{sga3:iv:scholium:3:3:3:1}{3.3.3.1});
注 \SGARefLink{sga3:iv:remark:3:3:3:2}{3.3.3.2} 也是新增的.}
\textit{若 \(R\) 是 \(X\) 中的有效等价关系, 则 \(p\) 是有效满态射
(\SGARefLink{sga3:iv:definition:1:3}{1.3}). 若 \(f:X\to Y\)
是有效满态射, 则 \(\mathscr R(f)\) 是 \(X\) 中的有效等价关系,
其一个商为 \(Y\). 因此, \(X\) 中的有效等价关系与 \(X\) 的有效商,
即以 \(X\) 为源的有效满态射的等价类之间, 有一个一一对应的``Galois''对应.}

\medskip
\noindent\SGARefTarget{sga3:iv:definition:3:3:3}\textbf{定义 3.3.3.}
\textit{\(X\) 中的等价关系 \(R\) 称为\emph{普遍有效的},
是指商 \(Y=X/R\) 存在, 且对每个 \(Y'\to Y\),
纤维积 \(X'=X\times_Y Y'\) 与 \(R'=R\times_Y Y'\) 都存在,
并且 \(R'\) 是 \(X'\) 在 \(Y'\) 上的纤维平方.
等价地, \(R\) 有效且 \(p:X\to X/R\) 为普遍有效满态射.}

\medskip
\noindent\SGARefTarget{sga3:iv:scholium:3:3:3:1}\textbf{附论 3.3.3.1.}\footnotemark[14]
\textit{与上面相同, \(X\) 中的普遍有效等价关系与
\(X\) 的普遍有效商之间存在一一对应.}

% Source page: Exp4 component-local p. 16; combined-reader p. 204; printed p. 192.
\medskip
\noindent\SGARefTarget{sga3:iv:remark:3:3:3:2}\textbf{注 3.3.3.2.}\footnotemark[14]
假设 \(\mathscr C\) 为 \(S\)-概形范畴, 并用 \(\mathbf A^1_S\)
表示 \(S\) 上的仿射直线. 设 \(R\subset X\times_S X\)
是普遍有效等价关系, \(p:X\to Y\) 为其商.
则对 \(Y\) 的每个开子集 \(U\),
\[
  \mathcal O(U)=\operatorname{Hom}_S(U,\mathbf A^1_S)
\]
是下列集合中满足 \(\varphi\circ\operatorname{pr}_1
=\varphi\circ\operatorname{pr}_2\) 的元素 \(\varphi\) 组成的集合:
\[
  \mathcal O(p^{-1}(U))
  =\operatorname{Hom}_S(p^{-1}(U),\mathbf A^1_S)
\]
特别地, 若 \(R\) 由群 \(H\) 在 \(X\) 上的自由右作用定义
(参见 \SGARefLink{sga3:iv:definition:3:2:1}{3.2.1}),
则 \(\mathcal O(U)\) 是所有满足下述条件的
\(\varphi\in\mathcal O(p^{-1}(U))\) 组成的集合:
\[
  \varphi(xh)=\varphi(x)
\]
对每个 \(S'\to S\), \(x\in X(S')\) 和 \(h\in H(S')\) 都成立.

\medskip
\noindent\SGARefTarget{sga3:iv:proposition:3:3:4}\textbf{命题 3.3.4.}
\textit{设 \(R\) 是 \(X\) 中的普遍有效等价关系.
设 \(f:X\to Z\) 是与 \(R\) 相容的态射,
因而经过某态射 \(g:X/R\to Z\) 分解. 以下条件等价:}
\begin{enumerate}
  \item[(i)]\SGARefTarget{sga3:iv:proposition:3:3:4:item:i:01}
  \textit{\(g\) 是单态射;}
  \item[(ii)]\SGARefTarget{sga3:iv:proposition:3:3:4:item:ii:01}
  \textit{\(R\) 是 \(f\) 定义的等价关系.}
\end{enumerate}

事实上, \SGARefLink{sga3:iv:proposition:3:3:4:item:i:01}{(i)}
显然蕴含 \SGARefLink{sga3:iv:proposition:3:3:4:item:ii:01}{(ii)},
逆命题则恰好是 \SGARefLink{sga3:iv:corollary:2:5}{2.5}.

\medskip
\noindent\SGARefTarget{sga3:iv:definition:3:3:5}\textbf{定义 3.3.5.}
\textit{设 \(H\) 是在 \(X\) 上自由作用的 \(\mathscr C\)-群.
若 \(H\) 的作用在 \(X\) 中定义的等价关系是有效的
(分别地, 普遍有效的), 则称 \(H\) \emph{有效地}作用,
或称 \(H\) 在 \(X\) 上的作用是\emph{有效的}
(分别地, \emph{普遍有效的}).}

\SGARefTarget{sga3:iv:subsection:3:4}\subsection*{3.4. \((M)\)-有效性}
\addcontentsline{toc}{subsection}{\SGARefLink{sga3:iv:subsection:3:4}{3.4}. (M)-有效性}

在实际应用中, 普遍有效满态射往往难以刻画.
然而, 我们经常已经掌握一些这种态射;
例如在概形论中, 忠实平坦拟紧态射就具有这一性质.
这就引出下面的讨论.

\medskip
\noindent\SGARefTarget{sga3:iv:para:3:4:1}\textbf{3.4.1.}
先对 \(\mathscr C\) 的一个态射族 \((M)\) 列出若干条件:
\begin{enumerate}
  \item[(a)]\SGARefTarget{sga3:iv:para:3:4:1:item:a:01}
  \((M)\) 在基变换下稳定: 属于 \((M)\) 的每个 \(u:T\to S\)
  都可平方 (参见 \SGARefLink{sga3:iv:definition:1:4:0}{1.4.0}),
  且对每个 \(S'\to S\), 基变换
  \[
    u':T\times_S S'\longrightarrow S'
  \]
  属于 \((M)\);
  \item[(b)]\SGARefTarget{sga3:iv:para:3:4:1:item:b:01}
  \((M)\) 的两个元素的复合属于 \((M)\);
  \item[(c)]\SGARefTarget{sga3:iv:para:3:4:1:item:c:01}
  每个同构都属于 \((M)\);
  \item[(d)]\SGARefTarget{sga3:iv:para:3:4:1:item:d:01}
  \((M)\) 的每个元素都是有效满态射.
\end{enumerate}

条件 \SGARefLink{sga3:iv:para:3:4:1:item:a:01}{(a)} 和
\SGARefLink{sga3:iv:para:3:4:1:item:b:01}{(b)} 蕴含:
\begin{enumerate}
  \item[(a$'$)]\SGARefTarget{sga3:iv:para:3:4:1:item:a:prime:01}
  \((M)\) 的两个元素的 Cartesian 积属于 \((M)\).
  更确切地说, 设 \(u:X\to Y\) 和 \(u':X'\to Y'\)
  是两个属于 \((M)\) 的 \(S\)-态射. 若 \(Y\times_S Y'\) 存在,
  则 \(X\times_S X'\) 存在且 \(u\times_Su'\) 属于 \((M)\).
\end{enumerate}

% Source page: Exp4 component-local p. 17; combined-reader p. 205; printed p. 193.
这可由下图得到
\[
\begin{array}{ccccc}
  &&Y'&\xleftarrow{u'}&X'\\
  &&\big\uparrow&&\big\uparrow\\[-2pt]
  X&\longleftarrow&X\times_S Y'&\longleftarrow&X\times_S X'\\
  {\scriptstyle u}\big\downarrow&&\big\downarrow&&
    \underset{\scriptstyle u\times_Su'}{\swarrow}\\[-2pt]
  Y&\longleftarrow&Y\times_S Y'&&
\end{array}
\]
类似地, \SGARefLink{sga3:iv:para:3:4:1:item:a:01}{(a)} 和
\SGARefLink{sga3:iv:para:3:4:1:item:d:01}{(d)} 蕴含:
\begin{enumerate}
  \item[(d$'$)]\SGARefTarget{sga3:iv:para:3:4:1:item:d:prime:01}
  \((M)\) 的每个元素都是普遍有效满态射.
\end{enumerate}

\medskip
\noindent\SGARefTarget{sga3:iv:para:3:4:2}\textbf{3.4.2.}
普遍有效满态射组成的族 \((M_0)\) 满足
\SGARefLink{sga3:iv:para:3:4:1}{3.4.1} 的条件
\SGARefLink{sga3:iv:para:3:4:1:item:a:01}{(a)}--\SGARefLink{sga3:iv:para:3:4:1:item:d:01}{(d)}.
事实上, \SGARefLink{sga3:iv:para:3:4:1:item:a:01}{(a)},
\SGARefLink{sga3:iv:para:3:4:1:item:c:01}{(c)} 和
\SGARefLink{sga3:iv:para:3:4:1:item:d:01}{(d)} 由定义成立,
\SGARefLink{sga3:iv:para:3:4:1:item:b:01}{(b)} 则由
\SGARefLink{sga3:iv:proposition:1:8}{1.8} 得到.
以下假设给定 \(\mathscr C\) 的一个满足这些条件的态射族 \((M)\);
所得结果特别适用于族 \((M_0)\).

\medskip
\noindent\SGARefTarget{sga3:iv:definition:3:4:3}\textbf{定义 3.4.3.}
\(X\) 中的等价关系 \(R\) 称为 \emph{\((M)\) 型的},
是指它可表且 \(p_1\in(M)\); 此时条件
\SGARefLink{sga3:iv:para:3:4:1:item:b:01}{(b)} 和
\SGARefLink{sga3:iv:para:3:4:1:item:c:01}{(c)} 蕴含 \(p_2\in(M)\).

关系 \(R\) 称为 \emph{\((M)\)-有效的}, 是指它有效,
且典范态射 \(X\to X/R\) 属于 \((M)\).

\(X\) 的商 \(Y\) 称为 \emph{\((M)\)-有效的},
是指典范态射 \(X\to Y\) 属于 \((M)\).

以下结果立即由 \SGARefLink{sga3:iv:definition:3:4:3}{定义} 得到.
\footnote{编者注: 为便于后文引用, 添加了编号
\SGARefLink{sga3:iv:proposition:3:4:3:1}{3.4.3.1},
并展开了 \SGARefLink{sga3:iv:proposition:3:4:3:1:item:i:01}{(i)} 的证明.}

\medskip
\noindent\SGARefTarget{sga3:iv:proposition:3:4:3:1}\textbf{命题 3.4.3.1.}
\begin{enumerate}
  \item[(i)]\SGARefTarget{sga3:iv:proposition:3:4:3:1:item:i:01}
  \textit{\((M)\)-有效等价关系是 \((M)\) 型的, 并且普遍有效.}
  \item[(ii)]\SGARefTarget{sga3:iv:proposition:3:4:3:1:item:ii:01}
  \textit{\((M)\)-有效商是普遍有效的
  (参见 \SGARefLink{sga3:iv:definition:3:3:3}{3.3.3}).}
  \item[(iii)]\SGARefTarget{sga3:iv:proposition:3:4:3:1:item:iii:01}
  \textit{映射 \(R\mapsto X/R\) 和 \(p\mapsto\mathscr R(p)\)
  在 \(X\) 中的 \((M)\)-有效等价关系与 \(X\) 的
  \((M)\)-有效商之间建立了一一对应.}
  \item[(iv)]\SGARefTarget{sga3:iv:proposition:3:4:3:1:item:iv:01}
  \textit{\((M_0)\)-有效等价于普遍有效.}
\end{enumerate}

证明 \SGARefLink{sga3:iv:proposition:3:4:3:1:item:i:01}{(i)}.
由于 \(R\) 是 \((M)\)-有效的, 有 Cartesian 方形
\[
  \begin{array}{ccc}
    R & \xrightarrow{p_2} & X\\
    {\scriptstyle p_1}\big\downarrow && \big\downarrow{\scriptstyle p}\\
    X & \xrightarrow{p} & X/R,
  \end{array}
\]
且 \(p\in(M)\). 由 \SGARefLink{sga3:iv:para:3:4:1:item:a:01}{3.4.1(a)},
\(p_1\) 和 \(p_2\) 属于 \((M)\), 故 \(R\) 是 \((M)\) 型的.

% Source page: Exp4 component-local p. 18; combined-reader p. 206; printed p. 194.
令 \(Y=X/R\), 设 \(Y'\to Y\) 为任意态射.
由 \SGARefLink{sga3:iv:para:3:4:1:item:a:01}{3.4.1(a)},
纤维积 \(X'=X\times_Y Y'\) 和 \(R'=R\times_Y Y'\) 存在,
且态射 \(X'\to Y'\) 和 \(p_i':R'\to X'\) 属于 \((M)\).
最后, 由于 \(R=X\times_Y X\), 由纤维积的结合性得到
\[
  R'=X\times_Y X\times_Y Y'=X'\times_{Y'}X'.
\]
这表明 \(R'\) 是 \((M)\)-有效的; 特别地, \(R\) 普遍有效.
这就证明了 \SGARefLink{sga3:iv:proposition:3:4:3:1:item:i:01}{(i)},
也证明了 \SGARefLink{sga3:iv:proposition:3:4:3:1:item:iv:01}{(iv)}.
结合定义 \SGARefLink{sga3:iv:definition:3:3:2}{3.3.2},
即可得到 \SGARefLink{sga3:iv:proposition:3:4:3:1:item:ii:01}{(ii)} 和
\SGARefLink{sga3:iv:proposition:3:4:3:1:item:iii:01}{(iii)}.

\medskip
\noindent\SGARefTarget{sga3:iv:para:3:4:4}\textbf{3.4.4.}
设 \(H\) 是结构态射属于 \((M)\) 的 \(S\)-群.
若 \(H\) 在 \(S\)-对象 \(X\) 上自由作用,
则它在 \(X\) 中定义一个 \((M)\) 型等价关系.
事实上, 由 \SGARefLink{sga3:iv:para:3:4:1:item:a:01}{(a)},
纤维积 \(H\times_S X\) 存在, 且
\(\operatorname{pr}_2:H\times_S X\to X\) 属于 \((M)\).
若 \(H\) 的作用在 \(X\) 中定义的等价关系是 \((M)\)-有效的,
则称这一作用为 \emph{\((M)\)-有效的}.

\medskip
\noindent\SGARefTarget{sga3:iv:proposition:3:4:5}\textbf{命题 3.4.5 (\((M)\)-有效性与基变换).}
\textit{设 \(R\) 是 \(S\) 上 \(X\) 中的一个 \((M)\)-有效等价关系,
并令 \(Y=X/R\). 设 \(S'\to S\) 为使 \(Y'=Y\times_S S'\)
存在的基变换. 则 \(X'=X\times_S S'\) 存在,
\(R'=R\times_S S'\) 是 \(S'\) 上 \(X'\) 中的
\((M)\)-有效等价关系, 并且}
\[
  X'/R'\simeq (X/R)'.
\]

事实上, 典范态射 \(X\to Y\) 和 \(R\to Y\) 属于 \((M)\);
因此由 \SGARefLink{sga3:iv:para:3:4:1:item:a:prime:01}{(a$'$)},
\(X'\) 和 \(R'\) 可表. 由纤维积的结合性, \(R'\) 是
\(X'\) 中由典范态射 \(X'\to Y'\) 定义的等价关系,
而该态射属于 \((M)\). 这就证明了
\SGARefLink{sga3:iv:proposition:3:4:5}{所述断言}.

\medskip
\noindent\SGARefTarget{sga3:iv:proposition:3:4:6}\textbf{命题 3.4.6 (\((M)\)-有效性与 Cartesian 积).}
\textit{设 \(R\) (分别地, \(R'\)) 是 \(S\) 上的
\(X\) (分别地, \(X'\)) 中的 \((M)\)-有效等价关系.
若 \((X/R)\times_S(X'/R')\) 存在, 则 \(X\times_SX'\) 存在,
\(R\times_SR'\) 是 \(S\) 上 \(X\times_SX'\) 中的
\((M)\)-有效等价关系, 并且}
\[
  (X\times_SX')/(R\times_SR')
  \simeq (X/R)\times_S(X'/R').
\]

令 \(Y=X/R\) 及 \(Y'=X'/R'\).
由 \SGARefLink{sga3:iv:para:3:4:1:item:a:prime:01}{(a$'$)},
纤维积 \(X\times_SX'\) 存在, 且典范态射
\[
  q:X\times_SX'\longrightarrow Y\times_SY'
\]
属于 \((M)\). 公式
\[
  (X\times_SX')\times_{(Y\times_SY')}(X\times_SX')
  \simeq (X\times_YX)\times_S(X'\times_{Y'}X')
\]
(所有未注明基底的纤维积均在 \(S\) 上取)
表明 \(R\times_SR'\) 是 \(X\times_SX'\) 中由 \(q\) 定义的等价关系,
证明完毕.

% Source page: Exp4 component-local p. 19; combined-reader p. 207; printed p. 195.
\medskip
\noindent\SGARefTarget{sga3:iv:para:3:4:7}\textbf{3.4.7.}
\footnote{编者注: \SGARefLink{sga3:iv:para:3:4:7}{本段} 为新增内容.}
假设 \(\mathscr C\) 有终对象 \(e\), \(f:G\to G'\)
是一个属于 \((M)\) 的 \(\mathscr C\)-群态射.
由 \SGARefLink{sga3:iv:para:3:4:1:item:a:01}{3.4.1(a)},
核 \(\operatorname{Ker}(f)\) 由 \(e\times_{G'}G\) 表示,
且态射 \(\operatorname{Ker}(f)\to e\) 属于 \((M)\).

另一方面, \(f\) 定义的等价关系, 与 \(\operatorname{Ker}(f)\)
在 \(G\) 上的作用 (比如右作用) 定义的等价关系相同:
它是由集合映射 \((g,h)\mapsto(g,gh)\) 定义的态射
\[
  G\times\operatorname{Ker}(f)\longrightarrow G\times G
\]
的像. 因此, 由附论 \SGARefLink{sga3:iv:scholium:3:3:2:1}{3.3.2.1}
得到 \SGARefLink{sga3:iv:corollary:3:4:7:1}{如下推论}.

\medskip
\noindent\SGARefTarget{sga3:iv:corollary:3:4:7:1}\textbf{推论 3.4.7.1.}
\textit{假设 \(\mathscr C\) 有终对象 \(e\), \(f:G\to G'\)
是一个属于 \((M)\) 的 \(\mathscr C\)-群态射.
则 \(\operatorname{Ker}(f)\) 在 \(G\) 上的作用是
\((M)\)-有效的, 且 \(G'\) 为商 \(G/\operatorname{Ker}(f)\).}

\SGARefTarget{sga3:iv:subsection:3:5}\subsection*{3.5. 用下降构造商}
\addcontentsline{toc}{subsection}{\SGARefLink{sga3:iv:subsection:3:5}{3.5}. 用下降构造商}

经常遇到的情形是, 不能直接构造商, 但作适当的基变换后就可以.
本小节给出适用于这一情形的有用判别法.

在 \SGARefLink{sga3:iv:definition:2:1}{2.1} 中,
我们定义了 \(S'\) 上的对象 \(X'\) 上相对于态射 \(S'\to S\) 的下降数据.

\medskip
\noindent\SGARefTarget{sga3:iv:definition:3:5:1}\textbf{定义 3.5.1.}
\textit{\(X'\) 上相对于 \(S'\to S\) 的下降数据称为\emph{有效的},
是指带有此下降数据的 \(X'\) 同构于 \(S\) 上的某个对象 \(X\)
在 \(S'\) 上的逆像, 后者带有其典范下降数据.}

若 \(S'\to S\) 为下降态射 (\SGARefLink{sga3:iv:definition:2:2}{2.2}),
则 \SGARefLink{sga3:iv:definition:3:5:1}{定义} 中的对象 \(X\)
在唯一同构的意义下唯一.
称 \(S'\to S\) 为有效下降态射 (\SGARefLink{sga3:iv:definition:2:2}{2.2}),
就是指它为下降态射, 且相对于它的每个下降数据都是有效的.

现在考虑 \(S\) 上的对象 \(X\) 中的等价关系 \(R\).
设 \(X'\), \(X''\), \(X'''\) 分别为 \(X\) 在
\[
  S',\qquad S''=S'\times_SS',\qquad
  S'''=S'\times_SS'\times_SS',
\]
上的逆像, \(R'\), \(R''\), \(R'''\) 为由 \(R\) 取逆像得到的等价关系.
假设 \(X'\) 中的 \(R'\) 是 \((M)\)-有效的,
并考虑作为 \(S'\) 上的对象的商 \(Y'=X'/R'\).
由 \SGARefLink{sga3:iv:proposition:3:4:5}{3.4.5},
它在 \(S''\) 上的两个逆像都同构于 \(X''/R''\).
因此, \(S'\)-对象 \(Y'\) 带有典范粘合数据.
同样应用 \(X'''/R'''\) 的唯一性, 可知这实际上是下降数据.
(此论证隐含假设所列的纤维积都存在. 特别地,
若 \(S'\to S\) 可平方, 例如它是下降态射, 则此条件成立.)

\medskip
\noindent\SGARefTarget{sga3:iv:proposition:3:5:2}\textbf{命题 3.5.2.}
\textit{设 \(R\) 是 \(S\) 上的对象 \(X\) 中的等价关系,
\(S'\to S\) 是普遍有效满态射. 假设每个在 \(S'\) 上的逆像属于
\((M)\) 的 \(S\)-态射本身也属于 \((M)\). 以下条件等价:}
\begin{enumerate}
  \item[(i)]\SGARefTarget{sga3:iv:proposition:3:5:2:item:i:01}
  \textit{\(R\) 在 \(X\) 中是 \((M)\)-有效的;}
  % Source page: Exp4 component-local p. 20; combined-reader p. 208; printed p. 196.
  \item[(ii)]\SGARefTarget{sga3:iv:proposition:3:5:2:item:ii:01}
  \textit{\(R'\) 在 \(X'\) 中是 \((M)\)-有效的,
  且 \(X'/R'\) 上的典范下降数据有效.}
\end{enumerate}
\textit{当这些条件成立时, 由 \(X'/R'\) 下降得到的对象典范同构于 \(X/R\).}

蕴含 \SGARefLink{sga3:iv:proposition:3:5:2:item:i:01}{(i)}\(\Rightarrow\)\SGARefLink{sga3:iv:proposition:3:5:2:item:ii:01}{(ii)}
仅是将 \SGARefLink{sga3:iv:para:3:4:4}{3.4.4} 改述为下降语言.
逆命题证明后, 由普遍有效满态射为下降态射这一事实
(\SGARefLink{sga3:iv:proposition:2:3}{2.3}),
下降得到的对象在唯一同构的意义下唯一,
从而得到 \SGARefLink{sga3:iv:proposition:3:5:2}{本命题} 的最后一个断言.

证明 \SGARefLink{sga3:iv:proposition:3:5:2:item:ii:01}{(ii)}\(\Rightarrow\)\SGARefLink{sga3:iv:proposition:3:5:2:item:i:01}{(i)}.
令 \(Y'=X'/R'\), 设 \(Y\) 为下降得到的对象.
典范态射 \(p':X'\to X'/R'=Y'\) 按构造与下降数据相容:
它在 \(S''\) 上的两个逆像都与典范态射 \(X''\to X''/R''\) 相同.
因此, 它是某个 \(S\)-态射 \(p:X\to Y\) 在 \(S'\) 上的逆像.
由于 \(p'\) 属于 \((M)\), 关于 \(S'\to S\) 的假设蕴含
\(p\) 也属于 \((M)\). 此外, 由于 \(p'\) 与等价关系 \(R'\) 相容,
再次利用普遍有效满态射为下降态射这一事实, 可知态射 \(p\) 与 \(R\) 相容.
因而有态射
\[
  R\longrightarrow X\times_YX.
\]
为证明 \(R\) 是 \((M)\)-有效的且 \(Y\) 同构于 \(X/R\),
只须证明此态射为同构. 从 \(S\) 到 \(S'\) 作基变换后,
由于 \(R'\) 有效, 它成为同构;
因而由与 \SGARefLink{sga3:iv:proposition:2:4}{2.4} 相同的理由,
它本身就是同构.

注意, 若 \((M)\) 是普遍有效满态射组成的族 \((M_0)\),
且 \(\mathscr C\) 有纤维积, 则
\SGARefLink{sga3:iv:proposition:3:5:2}{本命题} 的假设成立
(\SGARefLink{sga3:iv:corollary:1:10}{1.10}). 我们得到如下推论.

\medskip
\noindent\SGARefTarget{sga3:iv:corollary:3:5:3}\textbf{推论 3.5.3.}
\textit{假设 \(\mathscr C\) 有纤维积 (只须 \(S\) 上的纤维积存在即可).
设 \(R\) 是 \(S\) 上 \(X\) 中的等价关系,
\(S'\to S\) 是普遍有效满态射. 以下条件等价:}
\begin{enumerate}
  \item[(i)]\SGARefTarget{sga3:iv:corollary:3:5:3:item:i:01}
  \textit{\(R\) 在 \(X\) 中普遍有效;}
  \item[(ii)]\SGARefTarget{sga3:iv:corollary:3:5:3:item:ii:01}
  \textit{\(R'\) 在 \(X'\) 中普遍有效,
  且 \(X'/R'\) 上的典范下降数据有效.}
\end{enumerate}
\textit{当这些条件成立时, 由 \(X'/R'\) 下降得到的对象典范同构于 \(X/R\).}

% End of assigned scope. The Section 4 heading begins immediately here in
% the authority, on Exp4 component-local p. 20 (combined-reader p. 208).
% Translation statistics: source_lines=756; output_lines=774; inline_math=603; display_math=35; footnotes=7; pending_TODO=0.
% Review: complete assigned component; mathematical symbols, targets, references, tags, environments and footnotes checked; formula prose translated.
